\documentclass[11pt]{article}
\usepackage[margin=1in]{geometry}
\usepackage{amsmath,amssymb}
\usepackage{enumitem}
\setlength{\parindent}{0pt}
\setlength{\parskip}{6pt}
\begin{document}
\section*{STAT 412 HW07 --- Question 2}

\subsection*{Problem}

\begin{quote}
The proportion of adults living in a small town who are college graduates is estimated to be $p=0.4$. To test this hypothesis, a random sample of 15 adults is selected. If the number of college graduates in the sample is anywhere from 3 to 9, we shall not reject the null hypothesis that $p=0.4$; otherwise, we shall conclude that $p\ne0.4$.
\end{quote}

\subsection*{Solution}

Let $X$ be the number of college graduates in the sample. The rule says not to reject when
\[
3\le X\le 9,
\]
so the rejection region is
\[
X\le2 \quad\text{or}\quad X\ge10.
\]

\textbf{Part (a).} When $p=0.4$, $X\sim\operatorname{Binomial}(15,0.4)$, so
\[
\alpha=P(X\le2)+P(X\ge10)\approx 0.0609.
\]

\textbf{Part (b).} The type II error probability is the probability of failing to reject when the alternative value of $p$ is true:
\[
\beta=P(3\le X\le9).
\]

For $p=0.3$,
\[
\beta=P(3\le X\le9),\quad X\sim\operatorname{Binomial}(15,0.3)
      \approx 0.8695.
\]

For $p=0.5$,
\[
\beta=P(3\le X\le9),\quad X\sim\operatorname{Binomial}(15,0.5)
      \approx 0.8454.
\]

\textbf{Part (c).} A good test procedure should make both $\alpha$ and $\beta$ small. Here, $\alpha=0.0609$ is relatively small, but both values of $\beta$ are large. Therefore, this is not a good test procedure. Select the statement that says $\alpha$ is relatively small while at least one $\beta$ is relatively large.

\subsection*{Final Answer}

\fbox{\begin{minipage}{0.94\linewidth}
$\alpha=0.0609$. For $p=0.3$, $\beta=0.8695$; for $p=0.5$, $\beta=0.8454$. This is not a good test because $\alpha$ is relatively small but the $\beta$ values are large.
\end{minipage}}

\end{document}
